\section{General Duality Framework}
\label{sec:general_framework}

\paragraph{Subsystem equilibrium with finite coupling.}
Consider
\begin{equation}
  H_{\mathrm{tot}} = H_Q + H_X + H_{\mathrm{int}}, \qquad
  H_{\mathrm{int}} = g\,f\otimes B.
\end{equation}
Here \(g\) is a coupling amplitude, \(f\) is a subsystem operator, and \(B\) is a bath operator.
Define
\begin{equation}
  Z_X(\beta):=\Tr_X e^{-\beta H_X}, \qquad
  Z_Q(\beta):=\Tr_Q e^{-\beta H_Q},
\end{equation}
and the reduced equilibrium operator
\begin{equation}
  \bar{\rho}_Q(\beta):=\Tr_X e^{-\beta H_{\mathrm{tot}}},
\end{equation}
with normalized reduced state
\begin{equation}
  \rho_Q(\beta)=\frac{\bar{\rho}_Q(\beta)}{\Tr_Q \bar{\rho}_Q(\beta)}.
\end{equation}
The Hamiltonian of mean force is the exact generator satisfying
\begin{equation}
  e^{-\beta H_{\mathrm{MF}}(\beta)}
  := \frac{\bar{\rho}_Q(\beta)}{Z_X(\beta)}.
\end{equation}
No perturbative approximation is used in this definition
\cite{gelinThermodynamicsSubensembleCanonical2009a,hiltHamiltonianMeanForce2011,talknerColloquiumStatisticalMechanics2020,rivasStrongCouplingThermodynamics2020}.

\begin{definition}[Influence operator and thermal-influence partition]
\label{def:influence_operator}
Define \(S_\beta\) by
\begin{equation}
  e^{S_\beta}
  := e^{+\beta H_Q/2}\,\frac{\bar{\rho}_Q(\beta)}{Z_X(\beta)}\,e^{+\beta H_Q/2}.
\end{equation}
Equivalently,
\begin{equation}
  \frac{\bar{\rho}_Q(\beta)}{Z_X(\beta)}
  = e^{-\beta H_Q/2}\,e^{S_\beta}\,e^{-\beta H_Q/2}.
\end{equation}
With \(\rho_Q^{\mathrm{th}}=Z_Q^{-1}e^{-\beta H_Q}\), this becomes
\begin{equation}
  \rho_Q(\beta)
  = \frac{Z_Q(\beta)}{Z_{\mathrm{MF}}(\beta)}
  \big(\rho_Q^{\mathrm{th}}(\beta)\big)^{1/2}
  e^{S_\beta}
  \big(\rho_Q^{\mathrm{th}}(\beta)\big)^{1/2},
\end{equation}
where \(Z_{\mathrm{MF}}(\beta):=\Tr_Q e^{-\beta H_{\mathrm{MF}}(\beta)}\).
\end{definition}

\paragraph{Physical meaning of the partition.}
The bare Gibbs factor carries the reference equilibrium geometry set by \(H_Q\), while
\(e^{S_\beta}\) carries all coupling-induced dressing.
The duality maps introduced below act on representatives of this dressed family.
They do not change the underlying family itself. This thermal-plus-influence split is the
operator-level equilibrium analogue of influence-functional dressing
\cite{feynmanTheoryGeneralQuantum1963a,mccaulMeanForceHamiltoniansInfluence2026}.

\begin{definition}[Representative family]
\label{def:rep_family}
Let \(\mathcal{P}\) be a parameter manifold and \(\mathcal{O}\) a set of observables.
A representative family is a map
\begin{equation}
  \mathcal{R}: \mathcal{P}\to \mathcal{O}, \qquad
  p \mapsto \{O_k(p)\}_{k\in K},
\end{equation}
where each \(O_k\) is evaluated from one fixed analytic functional form, while \(p\) can be
re-labeled by discrete representative transformations.
\end{definition}

\paragraph{Crossover coordinate as influence strength.}
Assume there is a positive scalar \(\chi\) that controls crossover in the closed solution and define
\begin{equation}
  y:=\log\chi.
\end{equation}
The surface \(\chi=1\) (equivalently \(y=0\)) is the self-dual crossover manifold separating weak
(\(\chi<1\)) and strong (\(\chi>1\)) influence branches.
For one dominant channel,
\begin{equation}
  \chi(\beta,g,\vartheta)=g^2\chi_0(\beta,\vartheta),\qquad \chi_0>0,
\end{equation}
so the moving self-dual scale is
\begin{equation}
  g_\star(\beta,\vartheta)=\chi_0(\beta,\vartheta)^{-1/2}.
\end{equation}

\begin{definition}[Observable class used for regime transfer]
\label{def:admissible_class}
An observable \(Q\) is admissible if
\begin{equation}
  Q(\beta,g,\vartheta)=F\!\bigl(\beta,\vartheta;\phi_1(\chi),\ldots,\phi_n(\chi),u\bigr),
\end{equation}
where \(F\) is smooth in its arguments, each \(\phi_j\) has asymptotic expansions for
\(\chi\to0\) and \(\chi\to\infty\), and \(u\) is an oriented scalar (possibly absent) that flips
sign under orientation-gauge transformation.
\end{definition}

\paragraph{Why this class is sufficient.}
If a transformation acts only on \((\chi,u)\), then the same closed function \(F\) is reused.
This is the mechanism behind regime transfer: one does not alter the analytic form, only the
argument regime where it is evaluated.

\paragraph{Roadmap from framework to physics.}
The next sections introduce strong-weak representative reflection, orientation gauge, and their
composition into a duality geometry.
The spin-boson realization is treated only after this framework is complete.
